% s-253-p253-m002; original PDF 253
\Aphorism{s-253-p253-m002}{349}
Ich finde eine gewisse Hyperbolik des Ausdrucks und der Gebärde häufig bei pedantischen und mittelmässigen Naturen, ausserdem — oder wohl gerade deshalb — auch bei jungen Männern und bei den meisten Weibern. Die kraftvollen und stark bewegten Seelen sind dagegen nach aussenhin der Übertreibung öfter abhold als geneigt, sie sprechen leise, sachlich und nüancirt, beinahe druckfertig, sie bewegen sich nicht ohne eine massvolle Plastik — selten genug, dass man sie rhetorisch flackern und qualmen sieht. Formvollendung — ein Ideal, dem selten von Pinsel und Feder, seltener noch von Hand und Mund nachgetrachtet wird; bei den besten Dichtern findet sie sich kaum, in der besten Gesellschaft sucht man sie nicht einmal.
\par

